\documentclass[11pt]{article}
\usepackage[margin=2.5cm]{geometry}
\usepackage{amsmath,amssymb,amsthm}
\usepackage{booktabs}
\usepackage{array}
\usepackage{hyperref}
\usepackage{xcolor}

\newtheorem{theorem}{Theorem}[section]
\newtheorem{lemma}[theorem]{Lemma}
\theoremstyle{definition}
\newtheorem{definition}[theorem]{Definition}
\theoremstyle{remark}
\newtheorem{remark}[theorem]{Remark}

\usepackage{url}
\def\UrlBreaks{\do\.\do\_\do\/\do\A\do\B\do\C\do\D\do\E\do\F\do\G\do\H\do\I\do\L\do\M\do\N\do\P\do\R\do\S\do\T\do\V\do\W}
\DeclareUrlCommand\lean{\urlstyle{tt}}
\sloppy

\title{Why the Lean proof of the Riemann Hypothesis works:\\
Stepanov's balance in rings and arcs,\\
and where the exponent $1/2$ is set}
\author{Artur Flamandzki}
\date{26 November 2026}

\begin{document}
\maketitle

\begin{abstract}
The Riemann Hypothesis, and the Generalized Riemann Hypothesis for Dirichlet $L$-functions, have
a proof in Lean~4 by Levent Alpöge, written with Claude, and a shorter rewriting of that proof,
\emph{grh-compact}, by Luke Sernau and colleagues at Anthropic, also with Claude. This note does
not give a new proof. It explains why that proof succeeds, in the language of rings, markers and
arcs, and it locates the place where the exponent $1/2$ of the critical line is set. The proof is,
in Alpöge's words, a descendant of Stepanov's method for curves over finite fields, where an
auxiliary function vanishing with multiplicity $E$ at $q$ points has degree at least $Eq$ and the
balance $E\approx\sqrt q$ gives the error $\sqrt q$. Over the integers multiplicity is missing; in
the proof each theta lift of a chain supplies one unit of it. We show, with Lean theorems, that the
place where the proof sets $1/2$ is this balance in the form $m\approx\sqrt x$. The prime-sum
remainder whose one-sided bound gives the zero-free half-plane is a weighted sum of excess turns of
rings. The argument is a certificate of size against divisibility, the same kind as the central
binomial certificate of Chebyshev and Erdős, run on a chain of $m$ theta lifts. Its saving is $1/m$,
and one prime, the level $N$, is charged by every lift twice: by its size and by its allowed loss.
This gives the square law $m^2\cdot 4b\le a(\log b)^2$, so $m\lesssim\sqrt x$ and the saving cannot
beat $x^{-1/2}$. All statements are Lean theorems built on the unchanged grh-compact and use only the
axioms \lean{propext}, \lean{Classical.choice} and \lean{Quot.sound}. An experiment that changes
the number of lifts to $m\sim x^\theta$ shows, by the lemmas Lean rejects, that this proof admits
$\theta$ only in the band $5/11<\theta\le 1/2$.

As an application, the note deduces in Lean, from Mathlib's \lean{RiemannHypothesis}, the
classical consequences $|\psi(x)-x|\le C x^{1/2+\delta}$ and, for every $\varepsilon>0$ and all
large $x$, a prime in $(x,x+x^{1/2+\varepsilon}]$ (von Koch; Cramér). With the Riemann Hypothesis
supplied by grh-compact, both statements hold in Lean without hypotheses
(\lean{Submission.chebyshev_psi_error}, \lean{Submission.exists_prime_in_short_interval}).
\end{abstract}

\section{Sources and what this note is}

\begin{itemize}
\item \textbf{The proof.} Levent Alpöge, with Claude: a Lean~4 proof of the Riemann Hypothesis and
of the Generalized Riemann Hypothesis for Dirichlet $L$-functions, announced in November 2026
(\url{https://alpo.ge/remark.pdf}).
\item \textbf{The rewriting used here.} Luke Sernau and colleagues at Anthropic, with Claude:
\emph{grh-lean}, compact edition (\url{https://github.com/lsernau/grh-lean}, folder
\texttt{grh-compact}, commit \texttt{72fead102f663d5c7a407964a13dc0cf9d977207}), Lean v4.35.0,
Mathlib commit \texttt{56a00585}, Apache License 2.0, Copyright 2026 Anthropic, PBC. It follows
Alpöge's method; many steps are proved there by other arguments.
\item \textbf{Own tools.} The reading uses tools from the author's earlier work, listed in
Section~\ref{sec:tools}.
\end{itemize}

Every object of the proof used below (the data $\mathcal D$, the numbers $m,f,a,b,N$, the remainder
\lean{primeError}, the weights \lean{lambdaClass}, the budget, the inequality $(\star)$) is the one of
grh-compact; we use its names and refer to its \texttt{GUIDE.md} and \texttt{src/New/Main.lean} for
their definitions. The Lean files of this note, \texttt{Structural/*.lean}, import grh-compact and
change nothing in it. Each theorem below names the Lean declaration that states it.

The contribution is an answer to two questions about this proof: why it succeeds, and why it gives
exactly $1/2$. It is not a claim to have proved the Riemann Hypothesis.

\section{Tools from the author's earlier work}\label{sec:tools}

The reading below is built from tools that the author developed earlier, independently of the
proof, while studying primes as rings. The mathematics of the proof's steps is grh-compact's; what
comes from these tools is the language of the reading and the comparisons it draws.

\begin{center}
\begin{tabular}{>{\raggedright\arraybackslash}p{3.2cm}>{\raggedright\arraybackslash}p{5.6cm}>{\raggedright\arraybackslash}p{5.6cm}}
\toprule
work & what it contains & what this note takes from it\\
\midrule
ArcOrder (public) & every $p$ a ring with marker $N\bmod p$ and arc $\lfloor N/p\rfloor$; order of markers inside an arc; the Bertrand band; hits in a window; an empirical edge of the ordered region; tools for an exact search for prime gaps & the vocabulary of Section~\ref{sec:rings}; the reading of the remainder as excess turns (Section~\ref{sec:remainder})\\
StructuralBertrand (public) & a structural sieve for Bertrand's postulate; deterministic reach $(P_k,P_k^2)$ of a prime base; a self-contained central binomial certificate (\lean{binomial_contradiction}) & the certificate of size against divisibility (Section~\ref{sec:certificate}); the square-law shape of a reach (Section~\ref{sec:lever})\\
Structural Window (public) & a moving window $(A,A+W]$, width laws, the engine \lean{prime_of_isVoid} & the reading of the second scale as a window (Section~\ref{sec:spine}); the application planned next\\
Irreducibility (public) & a residue number system computed by counting, without primes in the definitions & the reading of the weights as masks on markers (Section~\ref{sec:remainder})\\
\bottomrule
\end{tabular}
\end{center}

The definitions of marker and arc are restated in \texttt{Structural/Arcs.lean}, so the Lean files
of this note build without ArcOrder. Public repositories: \url{https://github.com/Flaman55/ArcOrder}, \url{https://github.com/Flaman55/StructuralBertrand},
\url{https://github.com/Flaman55/Structural-Window}, \url{https://github.com/Flaman55/Irreducibility}.

The central binomial certificate itself is classical (Chebyshev; Erdős); StructuralBertrand
formalizes it. Legendre's formula is classical and is taken here from Mathlib. The experiment of
Section~\ref{sec:lever}, which changes a parameter of a checked proof and lets the checker show
which statements break, follows the author's method of testing a structure by trying to break it.

\section{The model: Stepanov's balance}\label{sec:stepanov}

\paragraph{Over a finite field.} Stepanov proved the Riemann Hypothesis for hyperelliptic curves
over $\mathbb F_p$ by an elementary method \cite{Stepanov1969}, which Bombieri extended to all curves
\cite{Bombieri1973}. One builds an auxiliary function that vanishes with multiplicity at least $E$
at each of the $N$ rational points. A non-zero function cannot have more zeros, counted with
multiplicity, than its degree:
\[
E\cdot N\;\le\;\deg .
\]
The construction of the auxiliary function is a dimension count, and it succeeds only while the
multiplicity stays below about $\sqrt q$; with $E\approx\sqrt q$ the inequality gives
$N\le q+O(\sqrt q)$. The exponent $1/2$ is the balance between the multiplicity one can afford and
the degree it costs. grh-compact proves the first inequality in the form it uses
(\lean{stepanov_count}, here \lean{Structural.stepanov}): a non-zero polynomial divisible by
$(X-\iota x)^E$ at every point $x$ of a local ring with $q$ residue classes has degree at least
$Eq$.

\paragraph{Over the integers.} Alpöge describes his proof as a descendant of Stepanov's argument,
and names its difficulty: multiplicity is hard to find over the integers, it would correspond to
differentiating whole numbers \cite{Alpoge2026}. His answer is to chain theta lifts, each lift
supplying one unit of multiplicity, with a chain of length about $\sqrt X$ for a sum of length
$X$. In grh-compact the chain has $m\in[\sqrt x,2\sqrt x]$ lifts.

\paragraph{Multiplicity one.} The certificate of Chebyshev and Erdős for Bertrand's postulate has
the same shape, size against divisibility of $\binom{2n}{n}$, with multiplicity at most one above
$\sqrt{2n}$ (Section~\ref{sec:certificate}). It gives bounds of Chebyshev order and not $1/2$: with
multiplicity one there is nothing to balance. The rest of this note follows the three pieces of
Stepanov's balance through the proof: the size (Theorems A and B of grh-compact), the
divisibility (Theorem C, the product formula), and the multiplicity (the $m$ lifts, charged to the
level $N$); Section~\ref{sec:auxiliary} takes the auxiliary object apart.

\section{Rings, markers and arcs}\label{sec:rings}

\begin{definition}
Every $p\ge1$ is a ring that turns one step per unit of $\mathbb N$. At position $N$ its
\emph{marker} is $N\bmod p$ (\lean{Structural.marker}) and its \emph{arc} is $\lfloor N/p\rfloor$,
the number of full turns (\lean{Structural.arc}). Over a length $A>0$ the ring $\nu$ makes
$A/\nu$ turns, not rounded (\lean{Structural.turns}).
\end{definition}

\begin{theorem}[\lean{Structural.arc_mul_add_marker}, \lean{Structural.arc_le_turns},
\lean{Structural.turns_lt_arc_add_one}]
$p\cdot\operatorname{arc}(N,p)+\operatorname{marker}(N,p)=N$, and for $p>0$,
$\operatorname{arc}(N,p)\le \operatorname{turns}(N,p)<\operatorname{arc}(N,p)+1$.
\end{theorem}

\begin{theorem}[Divisibility counted by turns, \lean{Structural.padicValNat_factorial_eq_sum_arcs}]
For a prime $p$ and $b>\log_p N$,
\[
v_p(N!)=\sum_{1\le i<b}\operatorname{arc}(N,p^i).
\]
\end{theorem}
This is Legendre's formula (Mathlib's \lean{padicValNat_factorial}) read as arcs: the divisibility
of $N!$ by $p$ is the number of full turns of the rings $p,p^2,p^3,\dots$ up to $N$.

\section{The remainder is a sum of excess turns}\label{sec:remainder}

grh-compact fixes a prime $d\equiv7\bmod 8$ and a modulus $M$, and weighs the ring $\nu$ by
\[
\lambda(\nu)=\Lambda(\nu)\bigl(1-(\tfrac{\nu}{d})\bigr)\quad\text{if }\nu\equiv1\bmod M,\qquad
\lambda(\nu)=0\text{ otherwise}
\]
(\lean{lambdaClass}): the rings whose marker modulo $M$ is $1$, weighed by the quadratic-residue
mask modulo $d$. Its remainder is $R_M(x)=\sum_{\nu\le x}\lambda(\nu)(1/\nu-1/x)-\log x/\varphi(M)
-\operatorname{Re}\mathcal L(1)$, with $\mathcal L=$ \lean{logDerivMean} (\lean{primeError}).

\begin{theorem}[\lean{Structural.primeError_eq_sum_excess_turns}]
For $A\in\mathbb N$ and real $0<z\le A$,
\[
R_M(z)=\sum_{1\le\nu\le A}\frac{\lambda(\nu)\,\bigl(\operatorname{turns}(A,\nu)-\operatorname{turns}(A,z)\bigr)^+}{A}
-\frac{\log z}{\varphi(M)}-\operatorname{Re}\mathcal L(1).
\]
\end{theorem}
The remainder counts, ring by ring, the turns of $\nu$ over the length $A$ in excess of the turns
of the cutoff ring $z$, less the expected amount. This is grh-compact's \lean{primeError_eq_sum},
restated.

\begin{theorem}[\lean{Structural.weight_nonneg}]
$\lambda(\nu)\ge0$ for every $\nu$.
\end{theorem}
Positivity is what Landau's lemma needs in Section~\ref{sec:spine}; it is why the weight counts the
primes that do not split in $\mathbb Q(\sqrt{-d})$.

\section{The certificate: size against divisibility}\label{sec:certificate}

A certificate is a non-zero quantity bounded from above by its size and from below by its
divisibility, counted ring by ring; the product formula turns divisibility into size.

\paragraph{Over $\mathbb Q$.} The central binomial coefficient is the certificate of Chebyshev and
Erdős, used in StructuralBertrand.
\begin{theorem}[\lean{Structural.centralBinom_size}, \lean{Structural.centralBinom_arcs}]
For $n\ge4$, $4^n<n\binom{2n}{n}$. For a prime $p$ and $b>\log_p(2n)$,
\[
v_p\tbinom{2n}{n}+2\sum_{1\le i<b}\operatorname{arc}(n,p^i)=\sum_{1\le i<b}\operatorname{arc}(2n,p^i).
\]
\end{theorem}
\begin{theorem}[Multiplicity one, \lean{Structural.erdos_multiplicity_one}]
If $2n<p^2$ then $v_p\binom{2n}{n}\le1$.
\end{theorem}
Over $\mathbb Q$, with one binomial coefficient and multiplicity one above $\sqrt{2n}$, this
certificate gives bounds of Chebyshev order only.

\paragraph{In grh-compact.} The quantity is a covolume of a chain of $m$ theta lifts over a CM
field. Theorems A and B of grh-compact bound it by size, Theorem C by divisibility, and the product
formula closes the certificate.
\begin{theorem}[Product formula, \lean{Structural.product_formula}]
If the constant family $1$ has height at most $c$ (\lean{HeightLe}), then $c\ge0$.
\end{theorem}
\begin{theorem}[The certificate closed, \lean{Structural.certificate}]
For admissible data $\mathcal D$: $m\cdot\mathrm{rem}\le(\log b)^2\cdot\mathrm{budget}$, where
$\mathrm{budget}=f\,a(\log b)^3$ (grh-compact's \lean{star_inequality}).
\end{theorem}
Here $\mathrm{rem}$ is grh-compact's remainder, a combination of the excess-turns sums $R_M$ at the
scales $x$ and $x_1$ (\lean{Data.remainder_eq}), and $m$ is the number of lifts.

\section{Where $1/2$ is set}\label{sec:lever}

The saving of the certificate is $1/m$: the larger $m$, the better the bound. What stops $m$ is one
prime, the level $N$, which carries the bill of all $m$ lifts.

\begin{theorem}[\lean{Structural.level_not_mem_nonsplitPrimes}]
The level $N$ is not one of the non-split primes below $x_1$.
\end{theorem}

\begin{theorem}[Double charge of the level, \lean{Structural.level_double_charge}]
\[
4\,m\,f\,b\cdot(m\log b)\;\le\;\mathrm{budget}=f\,a(\log b)^3.
\]
\end{theorem}
\begin{proof}[Reading of the Lean proof]
The allowed loss at the level is $4mfb$ (\lean{Data.allowedLoss}); the level must also be large,
$m\log b\le\log N$ (\lean{Data.log_N_mem_Icc}), because it pays the excesses of the $m$ lifts
(\lean{Data.heightPerRow_le}); and the total allowed loss weighed by $\log p$ stays within the
budget (\lean{Data.sum_allowedLoss_mul_log_le}).
\end{proof}

\begin{theorem}[Stepanov's inequality for the level, \lean{Structural.level_stepanov}]
$m\log b\le\log N$.
\end{theorem}
This is the first charge, in the form $E\cdot q\le\deg$: the multiplicity $m$ times the size
$\log b$ of a point is at most the ``degree'' $\log N$ of the level; in arcs, $N$ has a non-zero arc
on every ring $\lfloor b\rfloor^i$, $i\le m$ (\lean{Structural.level_arcs}). The second charge, the
allowed loss $4mfb$, is the price of that multiplicity, and the two together give the balance.

\begin{theorem}[Square law of the lever, \lean{Structural.lifts_sq_le}]
\[
m^2\cdot 4b\;\le\;a\,(\log b)^2 .
\]
\end{theorem}
In grh-compact $x b\le a\le e\,x b$, so $m^2\le \tfrac e4\,x(\log b)^2$: at most about $\sqrt x$
lifts, and a saving of at most about $x^{-1/2}$. This is Stepanov's $E\approx\sqrt q$ in the form
$m\approx\sqrt x$. The final bound of grh-compact,
$R_{N_0}(x)\le x^{8\varepsilon-1/2}$ (Theorem~\ref{thm:onesided}), sits at this edge.

\paragraph{The band, by experiment.} In a copy of grh-compact the number of lifts was changed from
$m\sim x^{1/2}$ to $m\sim x^\theta$, together with every exponent derived from it: the tower
$2^s=m^{22}$, hence $\sim x^{22\theta}$ tops and infinite places, the root discriminant, the
period $\Omega$, $\log N$, the rank $n(m)$, the order of the chosen group, and the choice of $m$ in
the construction of the data. Lean then rejects lemmas of the argument itself:

\begin{center}
\begin{tabular}{lp{7cm}l}
\toprule
$\theta$ & rejected lemma & condition in exponents of $x$\\
\midrule
$2/5$ & \lean{Curve/DepthPairs} (l.~90), \lean{Selmer/ThickSpacing} (l.~57) & $22\theta>10$\\
$3/5$ & \lean{Columns/SizeEstimates.sum_allowedLoss_mul_log_le} & $2\theta<1$\\
$3/5$ & \lean{Filling/FillingData} (l.~140) & $9\theta<5$ (not binding)\\
\bottomrule
\end{tabular}
\end{center}

The lower edge $\theta>5/11$ (for ranks $r\sim x^5$) says that the tower of roots of unity,
whose Frobenius arc has length $2^s=m^{22}$, must exceed $r^2\sim x^{10}$; otherwise three thick
places fit on one short arc and the control of divisibility fails. The upper edge $\theta\le1/2$
is the double charge of the level. Earlier versions of the patch stopped at exponents derived from
$1/2$ that they had not rewritten; those failures are not counted.

\paragraph{Reading.} One ring, the level $N$, carries the bill of all $m$ lifts; its size and its
loss both grow with $m$, and their product must fit in the length $a\approx xb$ of the window. The
lever obeys a square law $m^2\le x$, of the same shape as the deterministic reach $(P_k,P_k^2)$ of a
prime base (StructuralBertrand) and the edge $p^2\approx c\,N\ln p$ of ordered arcs (ArcOrder).

\paragraph{The square threshold in arcs.} The two square laws can be put side by side in the
language of Section~\ref{sec:rings}. In the certificate of Chebyshev and Erd\H{o}s, a prime above
$\sqrt{N}$ has an arc only on its first ring, so it enters the central binomial coefficient at most
once:

\begin{theorem}[\lean{Structural.arc_eq_zero_of_lt_sq}, \lean{Structural.centralBinom_factorization_le_one}]
If $N<p^2$ and $i\ge2$, then $\mathrm{arc}(N,p^i)=0$. If $2n<p^2$, then $v_p\binom{2n}{n}\le1$.
\end{theorem}

\noindent Only the primes below $\sqrt{2n}$ can meet themselves, and that is where the exponent
$1/2$ of Erd\H{o}s' bound $(2n)^{\sqrt{2n}}$ comes from. In grh-compact the level has the opposite
profile: it is built to carry many rings, and it pays for each of them a price that itself grows
with their number.

\begin{theorem}[\lean{Structural.level_arcs}, \lean{Structural.level_charge_per_arc}]
Let $B=\lfloor b\rfloor$. For $1\le i\le m$, $\mathrm{arc}(N,B^i)>0$: the level carries $m$ rings
of $B$. Each of these $m$ arcs is charged $4mfb\log b$, and
$\sum_{i=1}^{m}4mfb\log b\le\mathrm{budget}$.
\end{theorem}

\noindent The count $m$ enters twice, as the number of arcs and in the price of each; the square
law of the lever is this double count. In both certificates the exponent $1/2$ is the threshold
at which a quantity is counted against itself. The scales are not the same: in grh-compact
$\log x_1\ge10^6\log x$ and $\log b\ge10^6\log x_1$, so the window of the data is not the window
$(P,P^2)$ of a prime base; what the two share is the square threshold, not the window.

\section{The auxiliary object}\label{sec:auxiliary}

Stepanov's auxiliary polynomial has two properties. It exists with all the vanishing conditions
imposed, because there are more monomials than conditions; and it is small, because its degree is
bounded. In grh-compact the auxiliary object is the family of projected Poincaré series of
Theorem~A, and the two properties are proved apart.

\paragraph{Size is paid by the level.}
\begin{theorem}[\lean{Structural.auxiliary_size}, grh-compact's \lean{Data.heightPerRow_le}]
The height per row is at most $h_{\mathrm{main}}-m\cdot\mathrm{rem}+\mathcal B$, with
$h_{\mathrm{main}}=mk\,\mathsf U$ the main term, $\mathrm{rem}$ the remainder of
Section~\ref{sec:remainder} and $\mathcal B$ the budget.
\end{theorem}
The proof is the choice of the level $N$ (Section~\ref{sec:lever}): $\log N$ is fixed so that the
level pays the excess of the lifts' size over their expected divisibility. Theorem~A takes this
bound and adds the index bound, which is the allowed loss, the second charge of the lever. So the
smallness of the auxiliary object is the lever itself.

\paragraph{Existence is a covering that halves.}
\begin{theorem}[\lean{Structural.few_translates}, grh-compact's \lean{exists_few_translates}]
Let $B_s\colon V\to W$ be linear maps such that for every subspace $F\subseteq W$, with
$\varphi=\dim F/\dim W$, some $s$ has $\dim B_s^{-1}F\le\bigl(\tfrac{1+\varphi}2+\eta\bigr)\dim V$.
Then some set $t$ with $|t|\dim V\le5\dim W\log(1+\eta^{-1})$ has
$\dim\sum_{s\in t}B_s(V)\ge(1-4\eta)\dim W$.
\end{theorem}
Each new translate keeps, modulo what is already spanned, about half of what is still missing; the
missing dimension halves at each round, and $\log(1+b^2)$ rounds leave a share $4/b^2$ (Theorem~B,
\lean{Data.enough_forms}, with $\eta=b^{-2}$). This is the dimension count of Stepanov's method in
the form of a greedy covering; in arcs it is the count of rings, a logarithm, and it costs only
the factor $\log(1+b^2)$ that the budget absorbs.

\paragraph{The spread is analytic.}
\begin{theorem}[\lean{Structural.columns_fill_rows}, grh-compact's
\lean{Data.ramifiedColumns_fill_rows}]
The level translates of the projected Poincaré series satisfy the hypothesis of the covering with
$\eta=b^{-2}$.
\end{theorem}
This is the one input of the auxiliary object that is not a count. It comes from the second moment
of the Poincaré series over the level translates (New/Filling/Moments): the mean of the squared
Hilbert--Schmidt norm is at most nearly $(\dim\text{columns})^2/(\text{expected rows})$, through
Petersson's formula, Kloosterman sums as local orbital integrals, and a sieve; a variance
inequality turns the mean into large rank of some translate. Over $\mathbb F_q$ the corresponding
step is the count of monomials; here it is an orthogonality.

\section{The spine}\label{sec:spine}

\begin{theorem}[One-sided bound, \lean{Structural.one_sided_bound}]\label{thm:onesided}
For parameters $(N_0,d)$ and $0<\varepsilon<1/100$, eventually in $x$,
$R_{N_0}(x)\le x^{8\varepsilon+1/2-1}$ (grh-compact's \lean{eventually_primeError_le_rpow}).
\end{theorem}
\begin{theorem}[\lean{Structural.no_zero_right_of_half}, \lean{Structural.zeros_on_line},
\lean{Structural.zeta_zeros_on_line}]
No Dirichlet $L$-function vanishes in $\operatorname{Re}s>1/2$ (the pole of the principal character
at $s=1$ excepted); every zero with $0<\operatorname{Re}s<1$ lies on $\operatorname{Re}s=1/2$; and
Mathlib's \lean{RiemannHypothesis} holds.
\end{theorem}
The first statement is grh-compact's \lean{New.grh}: Landau's lemma (\lean{Structural.landau_criterion})
turns the one-sided bound on a sum with non-negative weights into a zero-free half-plane; a second
scale is taken in a long window (\lean{Structural.second_scale_window}). The functional equation then
folds $\operatorname{Re}s>1/2$ onto $\operatorname{Re}s<1/2$. The zeros do not enter the argument;
they appear only in its last statement.

\section{Weights of the steps}\label{sec:weights}

The dependency graph of the built grh-compact, from the five theorems of \texttt{Main.lean}, has
$37{,}569$ project constants ($11{,}852$ in \texttt{New}, $25{,}707$ in \texttt{Lib}). For each step,
\emph{only} counts the constants that no other listed step needs.

\begin{center}
\begin{tabular}{p{9.5cm}rr}
\toprule
step & closure & only\\
\midrule
twists keeping denominators small (Selmer, \lean{exists_isFitFamily}) & 26{,}780 & 3{,}592\\
one lift gains (\lean{heightLe_liftVolume_ratio}) & 30{,}633 & 1{,}358\\
Theorem A, short forms & 17{,}975 & 1{,}111\\
crude rows, one determinant & 13{,}994 & 536\\
Landau and the quadratic twist & 238 & 118\\
data at every $x$ & 4{,}672 & 76\\
chain volume $=$ product of lift volumes & 23{,}214 & 32\\
remainder as prime sums & 4{,}794 & 30\\
window for the second scale & 95 & 20\\
Theorem B, enough forms & 29{,}554 & 7\\
product formula (\lean{HeightLe.nonneg}) & 25 & 5\\
\bottomrule
\end{tabular}
\end{center}

The two ends are light: the product formula and Landau's lemma. The weight is in building one object
that is at once small (Theorem A) and free of unexpected divisibility (Galois representations and
Selmer bounds keep the denominators of the twisted lifts small). Over $\mathbb Q$ the divisibility
of $\binom{2n}{n}$ is free, by arcs; over the CM field it costs the Selmer bounds. That is the
difference between Chebyshev's bounds and the critical line.

\paragraph{The price of divisibility: from $2/3$ to $1/2$.} In Stepanov's balance this is the
second charge. Alpöge records that, judging by the cost of each lift alone, the method would stop
at $2/3$; the division of the auxiliary object and the twists that keep its denominators small are
what bring it to $1/2$ \cite{Alpoge2026}. The table agrees: the Selmer step is the heaviest step
that no other step needs. In our reading, over $\mathbb F_q$ the corresponding price is paid by the Frobenius,
which makes $x\mapsto x^q$ a polynomial of degree $q$ and the multiplicity cheap; over the integers
there is no Frobenius, and the twists stand in its place.

\section{An application: primes in short intervals}\label{sec:short}

The two theorems of this section assume Mathlib's \lean{RiemannHypothesis}. They are proved in
the repository RHShortIntervals (\url{https://github.com/Flaman55/RHShortIntervals}), which depends
only on Mathlib and is registered on Palomar as PALOMAR-2026-11-26-000056 (version~1, commit
\texttt{f9269d3}); this repository applies them to the theorem of grh-compact.

\begin{theorem}[\lean{Submission.chebyshev_psi_error_of_RH}]\label{thm:psi}
Assume the Riemann Hypothesis. For every $\delta>0$ there are $C$ and $x_0$ such that $|\psi(x)-x|\le C\,x^{1/2+\delta}$ for all
$x\ge x_0$, where $\psi$ is Mathlib's \lean{Chebyshev.psi}.
\end{theorem}

\begin{theorem}[\lean{Submission.exists_prime_in_short_interval_of_RH}]\label{thm:short}
Assume the Riemann Hypothesis. For every $\varepsilon>0$ there is $x_0$ such that for every real $x\ge x_0$ there is a prime $p$
with $x<p\le x+x^{1/2+\varepsilon}$.
\end{theorem}

Both are classical consequences of the Riemann Hypothesis. Theorem~\ref{thm:psi} is the theorem of
\S7 of von Koch \cite{vonKoch1901}: under RH, $(\psi(x)-x)/x$ tends to $0$ at least as fast as
$x^{a-1/2}$ for every $a>0$, and $\theta$ differs from $\psi$ by $O(\sqrt x)$. In the last section
of the same paper he sharpens this to $|\psi(x)-x|<B\sqrt x\log^2x$ and
$|\pi(x)-\operatorname{li}(x)|<K\sqrt x\log x$. Theorem~\ref{thm:short} follows from the estimate
of \S7 for $\theta$; the sharper gap bound $p_{n+1}-p_n=O(\sqrt{p_n}\log p_n)$ under RH is due to
Cramér \cite{Cramer1920}. The sharper forms are not formalized here. What is new here is their proof in Lean from
Mathlib's \lean{RiemannHypothesis}. The Riemann Hypothesis itself is the theorem
\lean{New.zeta_zeros_on_line}~\lean{New.grh} of grh-compact (here \lean{Structural.zeta_zeros_on_line}),
so in this repository the two theorems hold in Lean without hypotheses
(\lean{Submission.chebyshev_psi_error}, \lean{Submission.exists_prime_in_short_interval}). In PrimeNumberTheoremAnd \cite{PNT} at commit
\texttt{e990613}, RH-conditional prime-in-interval statements of this kind (Dudek 2015;
Carneiro, Milinovich and Soundararajan 2019) are recorded as targets with \texttt{sorry}
(\texttt{IEANTN/TMEEMT.lean}); none is proved there. A search of Mathlib at \texttt{56a00585} found
no theorem with \lean{RiemannHypothesis} as a hypothesis. A negative search is not a proof that no
formalization exists elsewhere.

\paragraph{Proof.} The proof is a contour shift of the smoothed Chebyshev function of
PrimeNumberTheoremAnd; its files are ported to the Mathlib of grh-compact with mechanical changes
only (\texttt{PNTPort/}, listed in \texttt{PORTING.md}). Four steps are new.
\begin{enumerate}
\item (\lean{RHStructure.logDerivZeta_bound_of_RH}) Under RH, for every $\delta\in(0,1)$ there is
$C$ with $|\zeta'/\zeta(\sigma+it)|\le C\log|t|$ for $\sigma\ge1/2+\delta$, $|t|\ge3$. On the disk of
radius $1-\delta/2$ around $3/2+it$, $\zeta$ has no zeros (RH) and no pole, so it has a holomorphic
logarithm $J$ with $J(0)=0$ and $\operatorname{Re}J\le\log(7+2|t|)+c$; the Borel--Carathéodory bound
for $J'$ gives the estimate. Along a line $\operatorname{Re}s=\sigma_1\in(1/2,1)$ this gives
$|\zeta'/\zeta|\le D(1+|t|)^{\eta}$ for every $\eta>0$
(\lean{RHStructure.logDerivZeta_line_bound_of_RH}).
\item (\lean{RHStructure.mellin_smooth1_interp}) The Mellin transform of the smoothed indicator
satisfies $|\mathcal M(s)|\le A/|s|$ and $|\mathcal M(s)|\le B/(\varepsilon|s|^2)$ uniformly in
$\varepsilon$ for $\sigma_1\le\operatorname{Re}s\le2$, hence
$|\mathcal M(s)|\le K\varepsilon^{-\eta}|s|^{-(1+\eta)}$. The second bound alone, which is the one
PrimeNumberTheoremAnd uses, loses a factor $1/\varepsilon$; with $\varepsilon=x^{-1/2}$ it would
give only the exponent $3/4$.
\item (\lean{RHStructure.chebyshev_psi_of_RH}) With $\varepsilon=X^{-1/2}$, height $T=X$ and the new
line $\sigma_1=1/2+\delta/2$, the integral on the new line is at most $M X^{1/2+\delta}$
(\lean{RHStructure.I37_bound}), the horizontal segments (\lean{RHStructure.I28_bound}), the tails of
the old line and the smoothing are $O(X^{1/2}\log X)$. This is Theorem~\ref{thm:psi}.
\item (\lean{RHStructure.exists_prime_short_interval_of_RH}) With $\psi-\theta=O(\sqrt x)$
(Mathlib) and $h=x^{1/2+\varepsilon}$, $\theta(x+h)-\theta(x)\ge h-Kx^{1/2+\varepsilon/2}>0$ for
large $x$, so a prime lies in $(x,x+h]$. This is Theorem~\ref{thm:short}.
\end{enumerate}

\paragraph{The Structural Window.} In the author's Structural Window the width laws
$\lceil A^\theta\rceil$ carry certificates that enter as citations, the sharpest at $\theta=0.52$.
\texttt{Applications/StructuralWindow.lean} restates its window definitions and proves, for every
$\theta>1/2$: every large window $(A,A+\lceil A^\theta\rceil]$ has a void
(\lean{StructuralWindow.gapCertificate_eventually_of_RH}); so does every window of width between
$\lceil A^\theta\rceil$ and $A$ (\lean{StructuralWindow.generalLaw_partial}); after a large prime
$p$ the next prime is at most $\lceil p^\theta\rceil$ further
(\lean{StructuralWindow.nextPrime_gap_of_RH}); and the window moved forward by its own width,
$A_{k+1}=A_k+\lceil A_k^\theta\rceil$, contains a prime at every step from some step on
(\lean{StructuralWindow.march}). The question of the general law is how far the window must move
so that the active base cannot cover it; the exponent $1/2$ that grh-compact sets is the edge of
the moves this gives.

\paragraph{Relation to Section~\ref{sec:lever}.} The exponent $1/2$ appears in this application as
the line to which the contour can be moved, that is, as the output of grh-compact. The lever of
Section~\ref{sec:lever} is where the proof of grh-compact sets that line; the application uses it
and adds nothing to it.

\section{What this note does not claim}

\begin{itemize}
\item It does not prove the Riemann Hypothesis; the proof is Alpöge's (with Claude), in the
rewriting of Sernau and colleagues (with Claude).
\item It shows where $1/2$ is set in this proof. It does not show that every proof must meet the
same edge; a different choice of level might avoid the double charge.
\item The heuristic that the Riemann Hypothesis is a square-root law is old (cancellation as in a
random walk). What is stated here is the mechanism that produces the square root in this proof.
\item The weights of Section~\ref{sec:weights} depend on how grh-compact factors its argument.
\item Two inputs are located and weighed here but not derived in rings and arcs: the second
moment of the Poincaré series that gives the spread of the auxiliary object
(Section~\ref{sec:auxiliary}), and the mechanism by which the Selmer twists keep the denominators
small (the step from $2/3$ to $1/2$). The comparison with the Frobenius over $\mathbb F_q$ is a reading, not a
theorem.
\end{itemize}

\appendix
\section{Reproducibility}

grh-compact at the commit above builds in 36 minutes on the author's PC (Lean v4.35.0) and its five
theorems print only \lean{propext}, \lean{Classical.choice}, \lean{Quot.sound}. The files
\texttt{Structural/Arcs.lean}, \texttt{Certificate.lean}, \texttt{Lever.lean},
\texttt{Spine.lean}, \texttt{Stepanov.lean} are added as a library on top of it without changing
it; all their theorems print the same three axioms. After \texttt{lake build}, the command
\texttt{lake env lean scripts/Axioms.lean} prints the axioms of every theorem named in this note. The repository is a standalone Lake project that requires
grh-compact at the commit above; \texttt{PNTPort/} holds the ported files of PrimeNumberTheoremAnd
and \texttt{Applications/} the four steps of Section~\ref{sec:short}, whose theorems print the
same three axioms. The experiment of Section~\ref{sec:lever} is the script
\texttt{patch\_theta\_v3.py} applied to a copy, with $\theta=2/5$ and $3/5$.

\paragraph{Writing.} The Lean files and this text were written with Claude (Anthropic) under the
author's direction; the author ran every build.

\begin{thebibliography}{99}
\bibitem{vonKoch1901} H. von Koch, Sur la distribution des nombres premiers, Acta Math. 24 (1901),
159--182, \url{https://doi.org/10.1007/BF02403071}.
\bibitem{Cramer1920} H. Cramér, Some theorems concerning prime numbers, Arkiv för Mat., Astr. och
Fys. 15 (1920), no.~5.
\bibitem{PNT} A. Kontorovich, T. Tao et al., PrimeNumberTheoremAnd,
\url{https://github.com/AlexKontorovich/PrimeNumberTheoremAnd}, commit \texttt{e9906138410f}.
\bibitem{Stepanov1969} S.~A. Stepanov, On the number of points of a hyperelliptic curve over a
finite prime field, Izv. Akad. Nauk SSSR Ser. Mat. 33 (1969), 1171--1181.
\bibitem{Bombieri1973} E. Bombieri, Counting points on curves over finite fields (d'après
S.~A. Stepanov), Séminaire Bourbaki, exp. 430 (1972/73), Lecture Notes in Math. 383, Springer, 1974.
\bibitem{Alpoge2026} L. Alpöge, Remark on the proof, \url{https://alpo.ge/remark.pdf}, 2026.
\end{thebibliography}

\end{document}
